% !TEX program = xelatex
% 理论力学 · 专题讲义 · 并矢（dyad / dyadic）
\documentclass[11pt]{ctexart}
%=====================================================================
% theme-violet.tex —— 讲义精讲 · 紫花主题
% 用法：\documentclass[11pt]{ctexart}
%       \input{assets/theme-violet}
%       \begin{document} ... \end{document}
% 只用 ctexart 自带能力 + xcolor/enumitem/fancyhdr/pifont/amsmath，
% 不依赖 tikz/tcolorbox，保证 MiKTeX 可直接编译。
%=====================================================================

\usepackage[a4paper,left=2.0cm,right=2.0cm,top=2.1cm,bottom=2.3cm,%
            headsep=12pt,footskip=22pt]{geometry}
\usepackage{amsmath,amssymb}
\usepackage{needspace}
\usepackage{xcolor}
\usepackage{array}
\usepackage{enumitem}
\usepackage{pifont}
\usepackage{fancyhdr}
\usepackage[hidelinks]{hyperref}

%---------------------------------------------------------------- 配色
\definecolor{violet}{RGB}{107,78,143}      % 主色：紫
\definecolor{violetdeep}{RGB}{63,44,92}    % 深紫：标题/强调
\definecolor{petal}{RGB}{183,159,208}      % 花瓣：装饰线、花饰
\definecolor{bandviolet}{RGB}{247,243,251} % 浅紫底：档案/提示框
\definecolor{bandcream}{RGB}{252,249,242}  % 米白底：易错/注意框
\definecolor{rulepurple}{RGB}{217,207,230} % 分隔线
\definecolor{inkgray}{RGB}{92,96,110}      % 次要文字

%---------------------------------------------------------------- 字体
% 有 Noto CJK 用 Noto，没有就退回 ctex 默认（Windows 下为宋体/黑体）
\IfFontExistsTF{Noto Serif CJK SC}{\setCJKmainfont{Noto Serif CJK SC}}{}
\IfFontExistsTF{Noto Sans CJK SC}{\setCJKsansfont{Noto Sans CJK SC}}{}

%---------------------------------------------------------------- 版式
\setlist{nosep,leftmargin=1.9em}
\linespread{1.28}
\setlength{\parskip}{3.5pt}
\setlength{\parindent}{0pt}
\setlength{\tabcolsep}{5pt}
\renewcommand{\arraystretch}{1.28}

%---------------------------------------------------------------- 花饰
% \ding{95} = 五瓣小花，\ding{96} = 星形花；都来自 pifont，必定存在
\newcommand{\orn}{\textcolor{petal}{\ding{95}}}          % 行内小标
\newcommand{\ornstar}{\textcolor{petal}{\ding{96}}}      % 星形花
\newcommand{\flowerrule}{%                                % 居中花饰分隔线
  \par\vspace{7pt}%
  \noindent\makebox[\linewidth]{\color{petal}%
    \rule{0.20\linewidth}{0.6pt}\hspace{8pt}{\small\ding{95}}\hspace{4pt}%
    \raisebox{1.1pt}{\tiny\ding{95}}\hspace{4pt}{\small\ding{95}}\hspace{8pt}%
    \rule{0.20\linewidth}{0.6pt}}%
  \par\vspace{7pt}}

%---------------------------------------------------------------- 标题与节头
\newcommand{\doctitle}[3]{%   #1 主标题  #2 副标题  #3 元信息（可用 \\ 分行）
  \begin{center}
    {\color{petal}\small \orn\enspace #2\enspace \ornstar}\\[7pt]
    {\LARGE\bfseries\color{violetdeep} #1}\\[10pt]
    {\color{petal}\rule{0.55\linewidth}{1pt}}\\[4pt]
    {\small\color{inkgray} #3}
  \end{center}\vspace{2pt}}

\newcommand{\hdr}[2]{%   #1 章节号（如 §0）  #2 标题
  \par\vspace{12pt}%
  \noindent\colorbox{violetdeep}{\color{white}\bfseries\small\;#1\;}\hspace{7pt}%
  {\color{violetdeep}\large\bfseries #2}\par\nobreak
  \vspace{3pt}{\color{petal}\hrule height 1.2pt}\vspace{6pt}}

\newcommand{\subhdr}[1]{%
  \par\vspace{6pt}%
  \noindent{\color{violet}\bfseries #1}\par\vspace{1pt}}

%---------------------------------------------------------------- 提示框
\newcommand{\callout}[1]{% 浅紫底：讲义档案、读法说明
  \par\medskip\noindent
  \fcolorbox{rulepurple}{bandviolet}{\parbox{\dimexpr\linewidth-2\fboxsep-2\fboxrule}{\vspace{2pt}#1\vspace{2pt}}}%
  \par\medskip}

\newcommand{\warnbox}[1]{% 米白底：易错、边界条件
  \par\medskip\noindent
  \fcolorbox{rulepurple}{bandcream}{\parbox{\dimexpr\linewidth-2\fboxsep-2\fboxrule}{\vspace{2pt}#1\vspace{2pt}}}%
  \par\medskip}

\newcommand{\gapitem}[2]{% #1 缺口类型（紫） #2 内容
  \par\vspace{2pt}\noindent
  \hangindent=1.4em\hangafter=1
  {\color{petal}\small\ding{96}}\enspace{\bfseries\color{violetdeep}#1}\enspace #2\par}

% 题干块（\qstem）：每道例题／作业的「详解」开头必须先复述题干与出处。
% 起因：2026-10-25 第五章讲义把七道作业的题面集中放在一页、详解另起一页，
% 读者翻到「第 4 题」时看不到条件，等于一道没有题干的题目。
% 规则：详解块内自带题干，读者不翻页就能读懂题在问什么。
\newcommand{\qstem}[1]{%
  \par\medskip\noindent
  \fcolorbox{violet}{bandviolet}{%
    \parbox{\dimexpr\linewidth-2\fboxsep-2\fboxrule}{%
      \vspace{2pt}%
      {\small\bfseries\color{violet}题干}\hspace{0.8em}#1%
      \vspace{2pt}}}%
  \par\medskip}

% v2: 保护题干与开头，长题可跨页，续页补上下文。
\newcommand{\qbreak}{\needspace{6\baselineskip}}
% 独立练习：ID、来源、题干、考点、答案标签（通常 ans:ID）。
\newcommand{\practice}[5]{%
  \qbreak\subhdr{练习 #1}\qstem{（#2）#3}%
  \textbf{考点}：#4\par
  详解见附录第~\pageref{#5}~页。\par\medskip}
% 附录开头：ID、来源、条件摘要；自动产生 ans:ID 标签。
\newcommand{\answerstart}[3]{%
  \qbreak\subhdr{#1\quad 详解}\label{ans:#1}\qstem{（#2）#3}}
\newcommand{\answercontinue}[2]{%
  \needspace{5\baselineskip}\subhdr{#1\quad 续解}%
  {\small\color{inkgray}本页接续：#2}\par\medskip}
% 延迟提供默认值，让旧讲义可在导言区保留自己的 \newcommand{\step}。
\AtBeginDocument{\providecommand{\step}[1]{\needspace{3\baselineskip}\par\medskip\textbf{#1}\par}}

% 前提旁注：只放在「需要理解的那一步」后面，给陈述 + 本题代入，不重证。
\newcommand{\aside}[1]{%
  \par\vspace{1pt}\noindent
  \hangindent=1.6em\hangafter=1
  {\small\color{inkgray}%
    \hspace*{0.6em}{\color{petal}$\triangleright$}\hspace{2pt}#1}%
  \par\vspace{2pt}}

%---------------------------------------------------------------- 表格
% 本机 MiKTeX 的 colortbl 与 array 版本不匹配（\rowcolor 会报错），
% 故表格统一用「上粗紫线 + 表头 + 中细线」的三线式，不依赖 colortbl。
\newcommand{\tabtop}{\par\vspace{3pt}{\color{violet}\hrule height 1.2pt}\vspace{1pt}}
\newcommand{\tabmid}{\par\vspace{1pt}{\color{petal}\hrule height 0.5pt}\vspace{1pt}}
\newcommand{\tabbot}{\par\vspace{3pt}{\color{rulepurple}\hrule height 0.5pt}\vspace{8pt}}
\newcommand{\thead}[1]{{\bfseries\color{violetdeep}#1}}

%---------------------------------------------------------------- 页眉页脚
\pagestyle{fancy}
\fancyhf{}
\renewcommand{\headrulewidth}{0.4pt}
\renewcommand{\headrule}{\hbox to\headwidth{\color{rulepurple}\leaders\hrule height \headrulewidth\hfill}}
\fancyhead[L]{\small\color{inkgray}\songti \doctag}
\fancyhead[R]{\small\color{inkgray}\songti \coursetag}
\fancyfoot[C]{\small\color{inkgray}\orn\enspace{\color{violet}\thepage}\enspace\orn}

% 每篇讲义在导言区设这两个值
\newcommand{\doctag}{}
\newcommand{\coursetag}{}
\newcommand{\settags}[2]{\renewcommand{\doctag}{#1}\renewcommand{\coursetag}{#2}}

%---------------------------------------------------------------- 数学记号
\newcommand{\vect}[1]{\boldsymbol{#1}}
\newcommand{\dd}{\,\mathrm{d}}
\newcommand{\ddx}{\frac{\mathrm{d}}{\mathrm{d}x}}
\addtolength{\jot}{2pt}



\setlength{\headheight}{22pt}
\linespread{1.22}
\raggedbottom
\settags{专题 · 并矢}{理论力学}

% 记号：矢量用粗斜体（与 Topic 8 讲义一致），二阶张量用粗正体
\newcommand{\ten}[1]{\mathbf{#1}}
\newcommand{\unitd}{\hat{\ten{1}}}
\newcommand{\step}[1]{\par\smallskip\noindent{\color{violetdeep}\bfseries #1}\hspace{0.35em}}

\let\orighdr\hdr
\renewcommand{\hdr}[2]{\par\needspace{5\baselineskip}\orighdr{#1}{#2}}
\let\origsubhdr\subhdr
\renewcommand{\subhdr}[1]{\par\needspace{3\baselineskip}\origsubhdr{#1}}

\setlength{\abovedisplayskip}{6pt plus 2pt minus 2pt}
\setlength{\belowdisplayskip}{6pt plus 2pt minus 2pt}

\begin{document}

\doctitle{并矢：从张量积到惯量张量}
{一个“矢量进、矢量出”的线性算子，如何写成九个分量}
{自编专题讲义\quad 2026-11-03\quad 记号与 Topic 8（刚体运动学）一致}

\callout{\textbf{这份材料解决什么问题。}你在 Topic 8 的 p9--p12 已经见过并矢：旋转算子写成 $S=\sum_k\vect{n}'_k\otimes\vect{n}_k$，单位张量写成 $\unitd=\vect{n}_k\otimes\vect{n}_k$。但那两页只给了“怎么用”，没给“为什么需要它、它的代数有多完整、它在力学里到底落在哪”。这份讲义按 \textbf{定义 $\to$ 代数 $\to$ 主轴 $\to$ 微积分 $\to$ 三处力学落点（惯量张量、欧拉方程、应力）} 铺开，最后用三道例题把工具跑一遍。}

\hdr{§0}{为什么需要并矢}

\subhdr{0.1\quad 三种“进—出”类型}

物理公式按“吃进什么、吐出什么”分三类：

\begin{center}
\tabtop
\begin{tabular}{@{}p{0.20\linewidth}@{\hspace{8pt}}p{0.33\linewidth}@{\hspace{8pt}}p{0.38\linewidth}@{}}
\thead{类型} & \thead{例子} & \thead{用什么表示} \\
\multicolumn{3}{@{}l@{}}{\color{petal}\rule{\linewidth}{0.5pt}}\\
矢量 $\to$ 标量 & 功 $W=\vect{F}\cdot\vect{r}$ & 一个矢量就够（点乘） \\
标量 $\to$ 标量 & $U=U(\vect{r})$ & 标量函数 \\
矢量 $\to$ 矢量 & 力矩 $\vect{M}=\vect{r}\times\vect{F}$；角动量 $\vect{L}=\ten{J}\cdot\vect{\omega}$ & 九个分量，即二阶张量 \\
\end{tabular}
\tabbot
\end{center}

第三类的问题是：给它一个矢量 $\vect{u}$，返回一个矢量 $\vect{v}$，而且必须线性——$\vect{v}=T(\vect{u})$，$T(\alpha\vect{u}+\beta\vect{w})=\alpha T(\vect{u})+\beta T(\vect{w})$。\textbf{一个矢量装不下九个自由度数}，所以必须换一种对象来表示 $T$。并矢就是这种对象最原始的写法：把两个矢量“并排”写在一起。

\subhdr{0.2\quad 三个落点}

本讲义后面会反复用到同一台机器，落点只有三处：

\begin{itemize}
\item \textbf{转动惯量}：$\vect{L}=\ten{J}\cdot\vect{\omega}$、$T=\frac12\vect{\omega}\cdot\ten{J}\cdot\vect{\omega}$。同一个 $\vect{\omega}$，绕不同方向转得到的 $\vect{L}$ 方向不同——这正是“$\ten{J}$ 不是标量”的直接证据。
\item \textbf{旋转算子}：$S=\sum_k\vect{n}'_k\otimes\vect{n}_k$，把固定系分量搬到动系（Topic 8 p10）。
\item \textbf{应力}：$\vect{t}(\vect{n})=\ten{\sigma}\cdot\vect{n}$。同一个点，取不同截面法向 $\vect{n}$，面力矢量不同。
\end{itemize}

\subhdr{0.3\quad 记号}

\begin{center}
\tabtop
\begin{tabular}{@{}ll@{}}
\thead{记号} & \thead{含义} \\
\multicolumn{2}{@{}l@{}}{\color{petal}\rule{\linewidth}{0.5pt}}\\
$\vect{a},\vect{b},\vect{n}$ & 矢量（粗斜体），分量 $a_i$ \\
$\vect{a}\otimes\vect{b}$ & 并矢：两矢量并列，\textbf{顺序有意义}，共九个分量 \\
$\ten{D},\ten{J},\ten{\sigma}$ & 二阶张量（粗正体），分量 $D_{ij}$ \\
$\unitd$ & 单位并矢 $\unitd=\vect{e}_1\otimes\vect{e}_1+\vect{e}_2\otimes\vect{e}_2+\vect{e}_3\otimes\vect{e}_3$ \\
$\ten{D}\cdot\vect{x}$，$\vect{x}\cdot\ten{D}$ & 右作用与左作用（二者一般\textbf{不相等}） \\
$\ten{A}:\ten{B}$ & 双点积 $A_{ij}B_{ij}$（结果是标量） \\
\end{tabular}
\tabbot
\end{center}

\warnbox{\textbf{三种“两矢量相乘”不一样，别混}\\[2pt]
$\vect a\cdot\vect b$：\textbf{内积}，结果是\textbf{标量}（一个数）；\\
$\vect a\times\vect b$：\textbf{叉积}，结果是\textbf{矢量}；\\
$\vect a\otimes\vect b$：\textbf{并矢}，结果是\textbf{二阶张量}（九个分量，本身不是矢量、也不是数）。\\[2pt]
本讲义统一写 $\otimes$，与 Topic 8 p9 的 $\vect n'_k\otimes\vect n_k$ 一致。有些教材把两个矢量并排写成 $ab$ 当并矢（Gibbs 的经典写法），但印刷上极容易与标量积 $ab$ 看混，本讲义不采用并置写法。}

\hdr{§1}{定义与三种读法}

\subhdr{1.1\quad 并矢就是“有序的矢量对”}

给定 $\vect{a},\vect{b}$，并矢 $\vect{a}\otimes\vect{b}$ 的定义方式不是说“它等于什么数”，而是规定\textbf{它怎样作用}：

\[
\boxed{\;(\vect{a}\otimes\vect{b})\cdot\vect{x}=\vect{a}\,(\vect{b}\cdot\vect{x})\;}
\qquad\text{右作用：先把 }\vect{x}\text{ 投影到 }\vect{b}\text{ 上，再沿 }\vect{a}\text{ 放大。}
\]

读法很重要：$\vect{b}\cdot\vect{x}$ 是标量，乘上 $\vect{a}$ 得矢量。所以 $\vect{a}\otimes\vect{b}$ 是一个“吃掉一个矢量、吐出一个矢量”的机器。

记号提醒：这里写的 $\otimes$ 是\textbf{张量积}，不是点乘——$\vect a\otimes\vect b$ 本身既不是标量也不是矢量，只有让它\textbf{作用}在别的矢量上，才得到一个矢量。

改写成分量：设 $\vect{a}=a_i\vect{e}_i$，$\vect{b}=b_j\vect{e}_j$，则

\[
\vect{a}\otimes\vect{b}=a_ib_j\,\vect{e}_i\otimes\vect{e}_j
\qquad\Longrightarrow\qquad
D_{ij}=a_ib_j ,
\]

九个分量排成矩阵正好是\textbf{列乘行}（外积）：$\ten{D}=\vect{a}\otimes\vect{b}^{\mathsf T}$。

\subhdr{1.2\quad 左右作用不同：这是并矢最容易错的地方}

左作用是另一条规则：
\[
\vect{x}\cdot(\vect{a}\otimes\vect{b})=(\vect{x}\cdot\vect{a})\,\vect{b}
\qquad\text{（先投影到 }\vect{a}\text{，再沿 }\vect{b}\text{ 放大）}.
\]

两者一般\textbf{不同}。取 $\vect{a}=\vect{e}_1$、$\vect{b}=\vect{e}_2$、$\vect{x}=\vect{e}_1+2\vect{e}_2+3\vect{e}_3$：

\[
(\vect{e}_1\otimes\vect{e}_2)\cdot\vect{x}=\vect{e}_1(\vect{e}_2\cdot\vect{x})=2\vect{e}_1,
\qquad
\vect{x}\cdot(\vect{e}_1\otimes\vect{e}_2)=(\vect{x}\cdot\vect{e}_1)\vect{e}_2=\vect{e}_2 .
\]

一个指向 $\vect{e}_1$，一个指向 $\vect{e}_2$，差别一眼可见。只有当 $\ten{D}$ 对称时（力学里的 $\ten{J}$、$\ten{\sigma}$ 都对称），左右作用才相同。

\begin{center}
\fcolorbox{violet}{bandviolet}{\parbox{0.94\linewidth}{\centering\vspace{3pt}
\small 一句话：\textbf{$\vect{a}\otimes\vect{b}$ 不对称，$\vect{a}\otimes\vect{b}+\vect{b}\otimes\vect{a}$ 才对称。}\\
\small 这也是“并矢不是两个矢量的乘积”的原因——它带着顺序信息。\vspace{3pt}}}
\end{center}

\subhdr{1.3\quad 与矩阵记号的字典}

\begin{center}
\tabtop
\begin{tabular}{@{}ll@{}}
\thead{并矢写法} & \thead{矩阵写法} \\
\multicolumn{2}{@{}l@{}}{\color{petal}\rule{\linewidth}{0.5pt}}\\
$\ten{D}=\vect{a}\otimes\vect{b}$ & $\ten D=\vect a\,\vect b^{\mathsf T}$（列乘行，即外积） \\
$\ten{D}\cdot\vect{x}$ & $\ten D\,\vect x$（矩阵乘列矢量） \\
$\vect{x}\cdot\ten{D}$ & $\vect x^{\mathsf T}\ten D$（行乘矩阵） \\
$\ten{A}\cdot\ten{B}$ & 矩阵乘法 $\ten A\ten B$ \\
$\ten{A}:\ten{B}$ & $\operatorname{tr}(\ten A^{\mathsf T}\ten B)$ \\
$\ten{D}^{\mathsf T}$ & 转置，$(\vect a\otimes\vect b)^{\mathsf T}=\vect b\otimes\vect a$ \\
$\unitd$ & 单位矩阵 $\ten 1$ \\
\end{tabular}
\tabbot
\end{center}

有了这张字典，线性代数里的一切（特征值、对角化、迹、行列式）都能直接搬到并矢上。\textbf{但并矢不是矩阵}：矩阵是“坐标下的账本”，并矢是不依赖坐标系的对象，换基时分量变、对象不变。

\subhdr{1.4\quad 两个天天用的特例}

\textbf{单位并矢。} $\unitd=\vect{e}_1\otimes\vect{e}_1+\vect{e}_2\otimes\vect{e}_2+\vect{e}_3\otimes\vect{e}_3$，满足
\[
\unitd\cdot\vect{x}=\vect{e}_k(\vect{e}_k\cdot\vect{x})=\vect{x}
\qquad\text{（对任意正交基都成立）}.
\]

\textbf{投影并矢。} 设 $\vect{n}$ 是单位矢量，令 $\ten{P}=\vect{n}\otimes\vect{n}$，则
\[
\ten{P}\cdot\vect{x}=\vect{n}(\vect{n}\cdot\vect{x})
\quad\text{（$\vect{x}$ 在 }\vect{n}\text{ 上的投影）},
\qquad
\ten{P}^2=\ten{P}\ \text{（再投一次不变）},
\]
而 $\unitd-\vect{n}\otimes\vect{n}$ 是“投影到垂直于 $\vect{n}$ 的平面”的算子。这两个 $\ten P$ 的迹分别是 $1$ 与 $2$——正对应“投影掉几个自由度”。

\hdr{§2}{代数运算}

\subhdr{2.1\quad 线性运算}

并矢的集合是九维线性空间：$\alpha\ten A+\beta\ten B$ 逐分量相加。任何二阶张量都能展开成
\[
\ten{D}=D_{ij}\,\vect{e}_i\otimes\vect{e}_j ,
\]
$\vect{e}_i\otimes\vect{e}_j$（九个）就是它的基。注意 $\vect{e}_1\otimes\vect{e}_2$ 与 $\vect{e}_2\otimes\vect{e}_1$ 是\textbf{两个不同的基元素}。

\subhdr{2.2\quad 转置}

\[
(\vect{a}\otimes\vect{b})^{\mathsf T}=\vect{b}\otimes\vect{a},
\qquad
(\ten A\cdot\ten B)^{\mathsf T}=\ten B^{\mathsf T}\cdot\ten A^{\mathsf T},
\qquad
(\ten D^{\mathsf T})_{ij}=D_{ji}.
\]
转置的判别式很好用：\textbf{对任意 $\vect x,\vect y$ 有 $(\ten D\cdot\vect x)\cdot\vect y=\vect x\cdot(\ten D^{\mathsf T}\cdot\vect y)$}。这一条在把“力 $\to$ 力矩”改写成“角速度 $\to$ 速度”时反复出现。

\subhdr{2.3\quad 对称与反对称分解}

\[
\ten D=\underbrace{\tfrac12(\ten D+\ten D^{\mathsf T})}_{\text{对称}}\;+\;
\underbrace{\tfrac12(\ten D-\ten D^{\mathsf T})}_{\text{反对称}} .
\]
力学里这两半各有明确身份：\textbf{对称半}给惯量、应力、应变；\textbf{反对称半}等价于一个\textbf{轴矢量}，正是叉乘。

具体地，若 $\ten A$ 反对称，则存在唯一的 $\vect a$ 使
\[
\ten A\cdot\vect x=\vect a\times\vect x
\qquad
\left(A_{ij}=-\varepsilon_{ijk}a_k,\quad a_k=-\tfrac12\varepsilon_{kij}A_{ij}\right).
\]
取 $\ten A=\vect r^{\times}$（Topic 8 p8 的对偶张量），就是 $\vect r^{\times}\vect x=\vect r\times\vect x$。$\vect a^{\times}$ 的矩阵是
\[
\vect a^{\times}=
\begin{pmatrix}
0&-a_3&a_2\\
a_3&0&-a_1\\
-a_2&a_1&0
\end{pmatrix}.
\]
\textbf{叉乘不是“另一种乘法”，而是一个反对称并矢在作用}——这句话能省掉很多“为什么叉乘不满足结合律”的困惑：不满足结合律的是叉乘这个运算，而反对称张量的合成（矩阵乘法）是满足的。

\subhdr{2.4\quad 迹、双点积、行列式}

\[
\operatorname{tr}\ten D=D_{ii},
\qquad
\ten A:\ten B=A_{ij}B_{ij}=\operatorname{tr}(\ten A^{\mathsf T}\cdot\ten B),
\qquad
\det\ten D=\varepsilon_{ijk}D_{1i}D_{2j}D_{3k}.
\]
两个常用恒等式：
\[
\ten A:\vect u\otimes\vect v=\vect u\cdot\ten A\cdot\vect v,
\qquad
\unitd:\ten D=\operatorname{tr}\ten D .
\]

\subhdr{2.5\quad 缩并层级表}

\begin{center}
\tabtop
\begin{tabular}{@{}lll@{}}
\thead{写法} & \thead{结果} & \thead{分量} \\
\multicolumn{3}{@{}l@{}}{\color{petal}\rule{\linewidth}{0.5pt}}\\
$\ten A\cdot\ten B$ & 二阶张量 & $A_{ik}B_{kj}$ \\
$\ten A\cdot\vect x$ & 矢量 & $A_{ij}x_j$ \\
$\ten A:\ten B$ & 标量 & $A_{ij}B_{ij}$ \\
$\vect a\otimes\vect b\cdot\vect c\otimes\vect d$ & 二阶张量 & $(\vect b\cdot\vect c)\,\vect a\otimes\vect d$ \\
$\vect a\cdot(\vect b\otimes\vect c)$ & 矢量 & $(\vect a\cdot\vect b)\vect c$ \\
\end{tabular}
\tabbot
\end{center}

最后一行那条 $\vect a\otimes\vect b\cdot\vect c\otimes\vect d=(\vect b\cdot\vect c)\vect a\otimes\vect d$ 是 Topic 8 p11 证明 $S^{\mathsf T}S=\unitd$ 的关键一步，值得单独记。

\hdr{§3}{主轴、谱分解、不变量}

\subhdr{3.1\quad 特征值问题}

若 $\ten D\cdot\vect n=\lambda\vect n$（$\vect n\ne\vect 0$），则 $\lambda$ 是特征值、$\vect n$ 是主轴方向。写成矩阵就是 $(\ten D-\lambda\unitd)\vect n=0$，有非零解的充要条件是
\[
\det(\ten D-\lambda\unitd)=0
\quad\Longleftrightarrow\quad
-\lambda^3+I_1\lambda^2-I_2\lambda+I_3=0 .
\]

\subhdr{3.2\quad 对称并矢的谱定理}

\textbf{对称}并矢（$\ten D=\ten D^{\mathsf T}$）必有三个实特征值 $\lambda_1,\lambda_2,\lambda_3$，且可挑三个互相正交的单位主轴 $\vect n_1,\vect n_2,\vect n_3$，使
\[
\ten D=\lambda_1\vect n_1\otimes\vect n_1+\lambda_2\vect n_2\otimes\vect n_2+\lambda_3\vect n_3\otimes\vect n_3 .
\]
这就是\textbf{谱分解}，也是“把张量搬到主轴上就变成三个数”的准确含义。力学里的惯量张量、应力张量都是对称的，所以这条定理在力学里\textbf{永远可用}。

\subhdr{3.3\quad 三个不变量}

换基时九个分量会变，但有三个组合不变，它们就是特征多项式的三个系数：

\[
I_1=\operatorname{tr}\ten D=\lambda_1+\lambda_2+\lambda_3,
\qquad
I_3=\det\ten D=\lambda_1\lambda_2\lambda_3,
\]
\[
I_2=\tfrac12\bigl[(\operatorname{tr}\ten D)^2-\operatorname{tr}\ten D^2\bigr]
=\lambda_1\lambda_2+\lambda_2\lambda_3+\lambda_3\lambda_1 .
\]
它们与坐标系无关：换基时九个分量变了，这三个数不变。做题时用来\textbf{验算}最好使。

\subhdr{3.4\quad 椭球的图像}

对对称正定的 $\ten D$，方程 $\vect x\cdot\ten D\cdot\vect x=1$ 是一个椭球，三个半轴长 $1/\sqrt{\lambda_k}$，方向沿主轴。\textbf{主轴就是椭球的轴}——惯量椭球、应力椭球都是这个图像。若某个 $\lambda=0$，椭球退化成柱面（这在“板”的惯量张量里会发生：板平面内的轴对应非零，法向对应另一个值，但绝不会出现零，因为刚体有质量分布；零出现于“质点系退化成线”等退化情形）。

\hdr{§4}{并矢的微积分}

\subhdr{4.1\quad 两个参照系下的时间导数（输运定理）}

固定系里看 $\ten D(t)$，与随刚体转动的动系里看，导数不同。设动系以角速度 $\vect\Omega$ 转动，则
\[
\left(\frac{\dd\ten D}{\dd t}\right)_{\text{固定}}
=
\left(\frac{\dd\ten D}{\dd t}\right)_{\text{动}}
+\vect\Omega\times\ten D-\ten D\times\vect\Omega .
\]
检验：取 $\ten D=\unitd$，左边为 $0$，右边 $\vect\Omega\times\unitd-\unitd\times\vect\Omega=\vect\Omega\times-\vect\Omega\times=0$，对上了。

这条式子是后面欧拉方程的\textbf{全部技术来源}：$\ten J$ 在动系（体坐标系）里是常张量，所以 $(\dd\ten J/\dd t)_{\text{动}}=0$，只剩两个叉乘项。

\subhdr{4.2\quad 散度、梯度与高斯定理}

两条最常用的公式与一条积分定理：

\[
(\nabla\cdot\ten D)_k=\partial_iD_{ik},
\qquad
\nabla\cdot(\vect a\otimes\vect b)=(\nabla\cdot\vect a)\,\vect b+(\vect a\cdot\nabla)\vect b,
\]
\[
\oint_S\ten D\cdot\vect n\dd S=\int_V\nabla\cdot\ten D\dd V .
\]
高斯定理的并矢版是连续介质力学的入口：把面力沿闭合曲面面积分换成体积分，就得到 $\nabla\cdot\ten\sigma+\vect f=\rho\vect a$。注意 $\nabla\cdot\ten\sigma$ 是\textbf{矢量}（对第二个指标缩并），而 $\nabla\ten\sigma$ 是三阶张量——两者别混。

\hdr{§5}{落点一：转动惯量张量}

\subhdr{5.1\quad 从 $\vect L=\int\vect r\times\vect v\dd m$ 推出来}

对刚体，$\vect v=\vect\omega\times\vect r$（基点取在瞬时不动点或质心），于是
\[
\vect L=\int\vect r\times(\vect\omega\times\vect r)\dd m
=\int\bigl[r^2\vect\omega-\vect r(\vect r\cdot\vect\omega)\bigr]\dd m .
\]
括号里对 $\vect\omega$ 是线性的，可以提出一个算子：
\[
\vect r\times(\vect\omega\times\vect r)=\bigl(r^2\unitd-\vect r\otimes\vect r\bigr)\cdot\vect\omega ,
\]
于是
\[
\boxed{\;\vect L=\ten J\cdot\vect\omega,\qquad
\ten J=\int\bigl(r^2\unitd-\vect r\otimes\vect r\bigr)\dd m\;}
\]
这就是转动惯量张量。\textbf{请注意它是并矢}：$r^2\unitd$ 是“各向同性部分”，$\vect r\otimes\vect r$ 是“沿 $\vect r$ 方向减掉的部分”。正因为有第二项，$\vect L$ 才一般\textbf{不平行}于 $\vect\omega$。

\subhdr{5.2\quad 分量与几何意义}

把 $\ten J$ 逐分量写出来：

\[
J_{ij}=\int\bigl(r^2\delta_{ij}-x_ix_j\bigr)\dd m
\qquad\Longrightarrow\qquad
\ten J=
\begin{pmatrix}
\int(y^2+z^2)\dd m & -\int xy\dd m & -\int xz\dd m\\[3pt]
-\int xy\dd m & \int(x^2+z^2)\dd m & -\int yz\dd m\\[3pt]
-\int xz\dd m & -\int yz\dd m & \int(x^2+y^2)\dd m
\end{pmatrix}.
\]
对角元是\textbf{转动惯量}（绕对应坐标轴），非对角元是\textbf{惯性积的负值}。$\ten J$ 对称、正定（只要质量不是全在一条线上），所以 §3 的谱定理全套可用。

\textbf{动能}：$T=\frac12\int v^2\dd m=\frac12\int(\vect\omega\times\vect r)^2\dd m=\frac12\vect\omega\cdot\ten J\cdot\vect\omega$。用并矢语言写，就是二次型。

\subhdr{5.3\quad 主轴 = 惯量主轴，主值 = 主转动惯量}

在主轴系里 $\ten J=\operatorname{diag}(A,B,C)$，于是
\[
\vect L=(A\omega_1,\;B\omega_2,\;C\omega_3),
\qquad
T=\tfrac12\bigl(A\omega_1^2+B\omega_2^2+C\omega_3^2\bigr).
\]
三个主转动惯量满足三角形不等式 $A+B\ge C$ 等（因为 $A+B-C=2\int z^2\dd m\ge0$）。这条不等式是很便宜的验算工具。

\subhdr{5.4\quad 两个必背的结果}

\textbf{长方体}（边长 $a,b,c$，质量 $m$，质心为原点，棱沿坐标轴）：
\[
\ten J=\frac{m}{12}\operatorname{diag}\bigl(b^2+c^2,\;a^2+c^2,\;a^2+b^2\bigr).
\]

\textbf{矩形薄板}（$a\times b$，厚度忽略，板面为 $xy$ 面）：
\[
\ten J=\frac{m}{12}\operatorname{diag}\bigl(b^2,\;a^2,\;a^2+b^2\bigr)
\qquad\Longrightarrow\qquad J_z=J_x+J_y ,
\]
最后这个等式就是\textbf{垂直轴定理}——它在并矢语言里只是“$z$ 方向没厚度，所以 $z$ 的惯量等于另两个之和”，一行就出来。

\textbf{平行轴定理的并矢版}（这个写法比标量版好用得多）：把原点从质心 $C$ 挪到 $O$，位移 $\vect R$，则
\[
\ten J_O=\ten J_C+m\bigl(R^2\unitd-\vect R\otimes\vect R\bigr).
\]
推导只用一条 $\int\vect\rho\dd m=0$（质心定义）。注意第二项是并矢，\textbf{不能}写成 $mR^2$ 一个标量——绕不同方向的轴平移，增量不同。

\hdr{§6}{落点二：欧拉方程与对称陀螺}

\subhdr{6.1\quad 体坐标系里的角动量定理}

固定系里 $\dd\vect L/\dd t=\vect M$。改用体坐标系（随刚体转，角速度 $\vect\omega$），由 §4.1 的输运公式：
\[
\left(\frac{\dd\vect L}{\dd t}\right)_{\text{固定}}
=\left(\frac{\dd\vect L}{\dd t}\right)_{\text{体}}+\vect\omega\times\vect L=\vect M .
\]
体坐标系里 $\ten J=\operatorname{diag}(A,B,C)$ 是常量，$\vect L=(A\omega_1,B\omega_2,C\omega_3)$，代入得\textbf{欧拉方程}
\[
A\dot\omega_1+(C-B)\omega_2\omega_3=M_1,
\quad
B\dot\omega_2+(A-C)\omega_3\omega_1=M_2,
\quad
C\dot\omega_3+(B-A)\omega_1\omega_2=M_3 .
\]
（符号的来源：$\vect\omega\times\vect L$ 的三个分量是 $(C-B)\omega_2\omega_3$、$(A-C)\omega_3\omega_1$、$(B-A)\omega_1\omega_2$。）

\subhdr{6.2\quad 对称陀螺的自由转动}

取 $A=B\ne C$、$\vect M=0$。第三条方程立刻给 $\dot\omega_3=0$，即 $\omega_3=\text{const}$；前两条变成
\[
\dot\omega_1=-\Omega\,\omega_2,\qquad \dot\omega_2=\Omega\,\omega_1,
\qquad
\Omega\equiv\frac{C-A}{A}\,\omega_3 .
\]
解是
\[
\omega_1=\omega_\perp\cos(\Omega t+\varphi_0),\qquad
\omega_2=\omega_\perp\sin(\Omega t+\varphi_0),
\]
即\textbf{角速度矢量绕体轴 $3$ 匀速转动}，横向分量大小不变、转速 $\Omega$。这就是“对称陀螺的自由进动”（这里进动的是 $\vect\omega$，不是陀螺本身）。

几何图像：$\vect L=(A\omega_1,A\omega_2,C\omega_3)$，与 $\vect\omega$ 都在由 $\vect e_3$ 与横向分量张成的\textbf{同一个平面}内，且这个平面绕 $\vect e_3$ 以 $\Omega$ 转动。三个矢量 $\vect\omega,\vect L,\vect e_3$ 始终共面——这是对称陀螺最有用的图像。

\begin{center}
\fcolorbox{violet}{bandviolet}{\parbox{0.94\linewidth}{\centering\vspace{3pt}
\small $C>A$（扁盘状，如地球）时 $\Omega>0$，$\vect\omega$ 与 $\vect L$ 绕 $\vect e_3$ 同向转；$C<A$（细长杆状）时反向。\vspace{3pt}}}
\end{center}

\subhdr{6.3\quad 为什么不能直接对 $\vect L$ 求导了事}

初学者常写 $\dd\vect L/\dd t=\ten J\cdot\dot{\vect\omega}$——这是\textbf{错的}，除非在固定系里 $\ten J$ 是常量。正确写法是
\[
\frac{\dd\vect L}{\dd t}=\ten J\cdot\dot{\vect\omega}+\dot{\ten J}\cdot\vect\omega
=\ten J\cdot\dot{\vect\omega}+\vect\omega\times\ten J\cdot\vect\omega ,
\]
第二项正是欧拉方程里那些 $\omega_i\omega_j$ 的来源。用并矢语言说：$\ten J$ 是“跟着体转”的张量，它的导数不为零。

\hdr{§7}{落点三：应力张量}

\subhdr{7.1\quad 面力是“法向的函数”}

在连续体里，过一点取一个截面，法向为 $\vect n$，截面上单位面积受力记 $\vect t(\vect n)$。它是 $\vect n$ 的\textbf{线性}函数（柯西应力定理），于是存在一个二阶张量 $\ten\sigma$ 使
\[
\vect t(\vect n)=\ten\sigma\cdot\vect n,
\qquad
\sigma_{ij}=\text{法向为 }\vect e_j\text{ 的面上、沿 }\vect e_i\text{ 方向的应力分量}.
\]
这就是并矢的第三个落点：\textbf{“九个应力分量”不是九个数随便摆，而是一个矢量到矢量的线性映射}。

\subhdr{7.2\quad 对称性来自角动量守恒}

对微元取力矩平衡（无体力矩），得到 $\sigma_{ij}=\sigma_{ji}$。所以应力张量也是\textbf{对称}的，§3 的谱定理同样适用：三个主应力 $\sigma_1,\sigma_2,\sigma_3$、三个互相垂直的主平面（剪应力为零的截面）。

\subhdr{7.3\quad 单轴拉伸：一个能口算的例子}

取 $\ten\sigma=\sigma_0\vect e_1\otimes\vect e_1$（只有 $\sigma_{11}=\sigma_0$）。主轴就是 $\vect e_1,\vect e_2,\vect e_3$，主应力 $\sigma_0,0,0$。

在 $xy$ 面内取 $45^\circ$ 截面，$\vect n=(\vect e_1+\vect e_2)/\sqrt2$：
\[
\vect t=\ten\sigma\cdot\vect n=\frac{\sigma_0}{\sqrt2}\vect e_1,
\]
法向分量 $=\vect t\cdot\vect n=\sigma_0/2$，剪向分量大小也是 $\sigma_0/2$。\textbf{最大剪应力出现在与主轴成 $45^\circ$ 的截面上，等于 $\sigma_0/2$}——这就是“拉伸试件沿 $45^\circ$ 断裂”的并矢解释。

\subhdr{7.4\quad 平衡方程}

对任意微元用高斯定理：$\oint_S\ten\sigma\cdot\vect n\dd S=\int_V\nabla\cdot\ten\sigma\dd V$，与体力 $\vect f$ 平衡得
\[
\nabla\cdot\ten\sigma+\vect f=0
\qquad\text{（静力学）},
\qquad
\nabla\cdot\ten\sigma+\vect f=\rho\vect a
\qquad\text{（动力学）}.
\]
写得开就是 $\partial_j\sigma_{ji}+f_i=0$——三个方程，正好对应三个方向的力平衡。

\hdr{§8}{三道例题}

\qstem{例 1\quad 设 $\ten D=4\vect e_1\otimes\vect e_1+\vect e_1\otimes\vect e_2+\vect e_2\otimes\vect e_1+4\vect e_2\otimes\vect e_2+2\vect e_3\otimes\vect e_3$。求 $\ten D\cdot\vect x$（$\vect x=\vect e_1+2\vect e_2+3\vect e_3$）、$\operatorname{tr}\ten D$、$\det\ten D$、$\ten D:\ten D$、特征值与主轴。}

\textbf{解。}矩阵是
\[
\ten D=\begin{pmatrix}4&1&0\\1&4&0\\0&0&2\end{pmatrix}.
\]
$\ten D\cdot\vect x=(4\cdot1+1\cdot2,\;1\cdot1+4\cdot2,\;2\cdot3)=(6,9,6)$；因为 $\ten D$ 对称，$\vect x\cdot\ten D$ 同值。

$\operatorname{tr}\ten D=4+4+2=10$；$\det\ten D=2(16-1)=30$；$\ten D:\ten D=16+1+1+16+4=38$。

特征方程 $-\lambda^3+10\lambda^2-31\lambda+30=0$，根为 $\lambda=5,3,2$。主轴：$\lambda=5$ 对应 $(\vect e_1+\vect e_2)/\sqrt2$；$\lambda=3$ 对应 $(-\vect e_1+\vect e_2)/\sqrt2$；$\lambda=2$ 对应 $\vect e_3$。谱分解
\[
\ten D=5\,\vect n_1\otimes\vect n_1+3\,\vect n_2\otimes\vect n_2+2\,\vect e_3\otimes\vect e_3 .
\]
验算：$5+3+2=10=\operatorname{tr}$，$5\cdot3\cdot2=30=\det$，都对。\checkmark

\qstem{例 2\quad 均匀长方体，质量 $m$，边长 $a=2,b=3,c=4$（米），质心在原点。求 $\ten J$；再把它压成薄板（$c\to0$，质量仍为 $m$），验证垂直轴定理。}

\textbf{解。}由 §5.4，
\[
\ten J=\frac{m}{12}\operatorname{diag}(b^2+c^2,\;a^2+c^2,\;a^2+b^2)
=\frac{m}{12}\operatorname{diag}(25,\;20,\;13).
\]
取 $m=5$ kg：$\ten J=\operatorname{diag}(125/12,\;100/12,\;65/12)\approx\operatorname{diag}(10.42,\;8.33,\;5.42)$ kg·m²。

压成薄板（$c=0$）后
\[
\ten J_{\text{板}}=\frac{m}{12}\operatorname{diag}(b^2,\;a^2,\;a^2+b^2)
=\frac{m}{12}\operatorname{diag}(9,\;4,\;13),
\]
于是 $J_z=13m/12$，$J_x+J_y=(9+4)m/12=13m/12$，二者相等，垂直轴定理成立。\checkmark

\qstem{例 3\quad 对称陀螺 $A=B=2$ kg·m²，$C=3$ kg·m²，自转 $\omega_3=10$ rad/s，无外力矩。求横向分量的进动角速度与周期。}

\textbf{解。}由 §6.2，
\[
\Omega=\frac{C-A}{A}\omega_3=\frac{3-2}{2}\times10=5\ \text{rad/s},
\qquad
T=\frac{2\pi}{\Omega}=\frac{2\pi}{5}\approx1.26\ \text{s}.
\]
即 $\vect\omega$ 的横向分量每 $1.26$ 秒绕体轴转一圈，$\omega_3$ 与 $|\vect\omega|$ 都不变。（数值积分欧拉方程的结果与之吻合到 $10^{-9}$。）

\hdr{§9}{易错点与速查}

\warnbox{\textbf{五个高频错误}\\[2pt]
1. \textbf{把 $\vect a\otimes\vect b$ 当成 $\vect a\cdot\vect b$}：前者是算子（九个分量），后者是标量。中间是 $\otimes$ 还是 $\cdot$，差一个层级；并矢表达式里最忌讳把 $\otimes$ 省成空位。\\
2. \textbf{左右作用混用}：$\ten D\cdot\vect x\ne\vect x\cdot\ten D$，除非 $\ten D$ 对称。力学里 $\ten J$、$\ten\sigma$ 对称，所以“碰巧”不错；换到旋转算子 $S$ 上立刻出错。\\
3. \textbf{直接写 $\dd\vect L/\dd t=\ten J\cdot\dot{\vect\omega}$}：只在 $\ten J$ 于该系为常量时成立；体坐标系里要补 $\vect\omega\times\ten J\cdot\vect\omega$。\\
4. \textbf{把惯量张量的非对角元符号写反}：$J_{12}=-\int xy\dd m$，负号不能丢；丢了以后主轴会算错。\\
5. \textbf{把 $\nabla\cdot\ten\sigma$ 当成标量}：它是对第二个指标缩并出来的\textbf{矢量}；写成 $\partial_i\sigma_{ij}$ 时注意哪个指标被缩并。}

\begin{center}
\begin{minipage}{\linewidth}
\tabtop
\begin{tabular}{@{}p{0.34\linewidth}p{0.56\linewidth}@{}}
\thead{要做的事} & \thead{用什么} \\
\multicolumn{2}{@{}l@{}}{\color{petal}\rule{\linewidth}{0.5pt}}\\
把“矢量 $\to$ 矢量”写成式子 & 并矢 $\ten D=D_{ij}\vect e_i\otimes\vect e_j$，作用规则 $(\vect a\otimes\vect b)\cdot\vect x=\vect a(\vect b\cdot\vect x)$\\[3pt]
判断是否对称 & 看 $D_{ij}=D_{ji}$；反对称部分等价于轴矢量 $\vect a$（叉乘）\\[3pt]
求主轴与主值 & $\det(\ten D-\lambda\unitd)=0$；对称时主轴正交，$\ten D=\sum\lambda_k\vect n_k\otimes\vect n_k$\\[3pt]
验算 & $I_1=\operatorname{tr}$、$I_3=\det$；惯量张量还有 $A+B\ge C$\\[3pt]
转动惯量 & $\ten J=\int(r^2\unitd-\vect r\otimes\vect r)\dd m$；平行轴 $\ten J_O=\ten J_C+m(R^2\unitd-\vect R\otimes\vect R)$\\[3pt]
角动量与动能 & $\vect L=\ten J\cdot\vect\omega$，$T=\frac12\vect\omega\cdot\ten J\cdot\vect\omega$\\[3pt]
体坐标系动力学 & $\dd\vect L/\dd t|_{\text{体}}+\vect\omega\times\vect L=\vect M$，即欧拉方程\\[3pt]
面力 & $\vect t(\vect n)=\ten\sigma\cdot\vect n$；$\sigma_{ij}=\sigma_{ji}$；主应力 = 特征值\\
\end{tabular}
\tabbot
\end{minipage}
\end{center}

\hdr{附录}{指标记法速查}

\begin{center}
\begin{minipage}{\linewidth}
\small
\renewcommand{\arraystretch}{1.05}
\tabtop
\begin{tabular}{@{}ll@{}}
\thead{对象} & \thead{指标写法} \\
\multicolumn{2}{@{}l@{}}{\color{petal}\rule{\linewidth}{0.5pt}}\\
并矢 & $D_{ij}=a_ib_j$ \\
右作用 & $(\ten D\cdot\vect x)_i=D_{ij}x_j$ \\
左作用 & $(\vect x\cdot\ten D)_j=x_iD_{ij}$ \\
张量乘张量 & $(\ten A\cdot\ten B)_{ij}=A_{ik}B_{kj}$ \\
双点积 & $\ten A:\ten B=A_{ij}B_{ij}$ \\
迹 & $\operatorname{tr}\ten D=D_{ii}$ \\
叉乘 & $(\vect a\times\vect b)_i=\varepsilon_{ijk}a_jb_k$ \\
对偶张量 & $(\vect a^{\times})_{ij}=-\varepsilon_{ijk}a_k$ \\
转动惯量 & $J_{ij}=\int(r^2\delta_{ij}-x_ix_j)\dd m$ \\
散度 & $(\nabla\cdot\ten D)_k=\partial_iD_{ik}$ \\
\end{tabular}
\tabbot
\end{minipage}
\end{center}

\aside{爱因斯坦求和约定：同一项里重复出现两次的指标自动求和。这份讲义只在附录与少数推导里用它，正文一律写求和号，避免“到底缩并了哪个指标”看错。}

\subhdr{延伸\quad 这套语言接下来怎么用}

\begin{itemize}
\item \textbf{刚体动力学的完整版}：欧拉角 + 欧拉方程 + 对称陀螺的进动与章动。Topic 8 的 Part 2 就是这条线；另外 Goldstein 第 4--5 章、Landau《力学》§25--§35 都用并矢记号写。
\item \textbf{二阶张量的一般理论}：谱定理、极分解、张量函数。任何一本连续介质力学教材的入门章都会系统讲，那时你会发现这份材料里的每条规则都还是原样。
\item \textbf{应力与应变}：弹性力学的标准语言就是并矢/张量；Landau《弹性理论》§1--§2 用九分量写应力张量，起点与 §7 完全一样。
\end{itemize}

\subhdr{自测五问}

\begin{enumerate}[label=(\arabic*)]
\item $\vect a\otimes\vect b\cdot\vect c$ 与 $\vect a(\vect b\cdot\vect c)$ 是否相同？那 $\vect a\otimes\vect b\cdot\vect c$ 与 $(\vect a\cdot\vect b)\vect c$ 呢？
\item 反对称张量 $\ten A$ 对应的轴矢量 $\vect a$ 满足什么关系？$\operatorname{tr}\vect a^{\times}$ 与 $\det\vect a^{\times}$ 各是多少？
\item $J_{12}=-\int xy\dd m$ 里的负号从哪来？把它丢掉以后，谱定理还成立吗？算出的主转动惯量还对吗？
\item 对称陀螺 $A=B\ne C$，若初始时 $\omega_3=0$、横向分量非零，运动是什么样？此时 §6.2 的 $\Omega$ 还有意义吗？
\item 单轴拉伸 $\ten\sigma=\sigma_0\vect e_1\otimes\vect e_1$，最大剪应力多大、出现在哪个截面？
\end{enumerate}

\aside{答案：(1) 相同，都等于“$\vect b\cdot\vect c$ 乘 $\vect a$”；但 $\vect a\otimes\vect b\cdot\vect c\ne(\vect a\cdot\vect b)\vect c$——后者是“$\vect a\cdot\vect b$ 乘 $\vect c$”，方向完全不同。(2) $A_{ij}=-\varepsilon_{ijk}a_k$，$a_k=-\frac12\varepsilon_{kij}A_{ij}$；迹与行列式都是 $0$（三维反对称矩阵必奇异，有一个零特征值）。(3) 负号来自 $J_{ij}=\int(r^2\delta_{ij}-x_ix_j)\dd m$ 的第二项；丢掉后矩阵仍对称，谱定理照样成立，但那是另一个张量，主轴与主值都不再是真实的惯量主轴与主转动惯量，$T=\frac12\vect\omega\cdot\ten J\cdot\vect\omega$ 会算错。(4) $\omega_1,\omega_2$ 恒定、$\omega_3=0$，即绕主轴的\textbf{永久转动}；此时 $\Omega=(C-A)\omega_3/A=0$，“进动”退化为静止。(5) $\sigma_0/2$，出现在与 $\vect e_1$ 成 $45^\circ$ 的截面上（§7.3）。}

\flowerrule

{\small\color{inkgray} 文中数值结果均经符号/数值复核：长方体与薄板的惯量张量由三重积分算出，例 1 的特征值与迹、行列式、双点积一致，对称陀螺的进动角速度与欧拉方程数值积分吻合。}

\end{document}
